
\documentclass{article}
\usepackage{colortbl}
\usepackage{makecell}
\usepackage{multirow}
\usepackage{supertabular}

\begin{document}

\newcounter{utterance}

\centering \large Interaction Transcript for game `matchit\_ascii', experiment `different\_grid', episode 2 with NEW\_Reflect\_Agent.
\vspace{24pt}

{ \footnotesize  \setcounter{utterance}{1}
\setlength{\tabcolsep}{0pt}
\begin{supertabular}{c@{$\;$}|p{.15\linewidth}@{}p{.15\linewidth}p{.15\linewidth}p{.15\linewidth}p{.15\linewidth}p{.15\linewidth}}
    \# & $\;$A & \multicolumn{4}{c}{Game Master} & $\;\:$B\\
    \hline

    \theutterance \stepcounter{utterance}  
    & & \multicolumn{4}{p{0.6\linewidth}}{
        \cellcolor[rgb]{0.9,0.9,0.9}{
            \makecell[{{p{\linewidth}}}]{
                \texttt{\tiny{[P1$\langle$GM]}}
                \texttt{You are participating in a collaborative guessing game.} \\
\texttt{The goal is to find out whether this grid and another picture only I can see, are the same.} \\
\texttt{Please describe your grid first. Then, I will provide my description and we can ask each other questions about the grids to figure out whether they are the same. Now start your short grid description with "DESCRIPTION:" followed by the description. Do not add anything else. This is your grid:} \\
\texttt{X X X X X} \\
\texttt{X X X X X} \\
\texttt{▢ ▢ X ▢ ▢} \\
\texttt{X X X X X} \\
\texttt{X X X X X} \\
            }
        }
    }
    & \\ \\

    \theutterance \stepcounter{utterance}  
    & \multicolumn{4}{p{0.45\linewidth}}{
        \cellcolor[rgb]{0.95,0.95,0.95}{
            \makecell[{{p{\linewidth}}}]{
                \texttt{\tiny{[P1$|$internal]}}
                \texttt{Reflections from previous episodes:} \\
\texttt{{-}} \\
\texttt{ Failed. I described the grid as a cross shape, but the user's unseen grid was different. My description was accurate but not detailed enough to catch subtle differences. Next time, I would specify exact positions of symbols and count them precisely to enable better comparison.} \\
\texttt{{-}} \\
\texttt{ Failed. I gave a clear description of my grid (border of X's with interior spaces). The user then described a different pattern (central cross shape). Our grids were not the same, which we determined after exchanging descriptions. My initial description was accurate and sufficient for comparison. Next time, I would follow the same approach: provide a concise, accurate initial description to enable quick identification of match or mismatch.} \\
\texttt{{-}} \\
\texttt{ Failed. I described the grid but did not ask clarifying questions about the other player's hidden grid to compare effectively. Next time, after giving my description, I should actively ask specific questions about their grid's layout to identify matches or mismatches.} \\
            }
        }
    }
    & & \\ \\

    \theutterance \stepcounter{utterance}  
    & \multicolumn{4}{p{0.45\linewidth}}{
        \cellcolor[rgb]{0.95,0.95,0.95}{
            \makecell[{{p{\linewidth}}}]{
                \texttt{\tiny{[P1$|$internal]}}
                \texttt{{[}attempt 1 – valid{]}} \\
\texttt{DESCRIPTION: A 5x5 grid with X's in all positions except for the center row (row 3), where only the middle cell is X and the other four cells are empty (▢). The first, second, fourth, and fifth rows are entirely filled with X's.} \\
            }
        }
    }
    & & \\ \\

    \theutterance \stepcounter{utterance}  
    & \multicolumn{4}{p{0.6\linewidth}}{
        \cellcolor[rgb]{0.9,0.9,0.9}{
            \makecell[{{p{\linewidth}}}]{
                \texttt{\tiny{[P1$\rangle$GM]}}
                \texttt{DESCRIPTION: A 5x5 grid with X's in all positions except for the center row (row 3), where only the middle cell is X and the other four cells are empty (▢). The first, second, fourth, and fifth rows are entirely filled with X's.} \\
            }
        }
    }
    & & \\ \\

    \theutterance \stepcounter{utterance}  
    & & & \multicolumn{2}{p{0.3\linewidth}}{
        \cellcolor[rgb]{0.9,0.9,0.9}{
            \makecell[{{p{\linewidth}}}]{
                \texttt{\tiny{[GM$|$GM]}}
                \texttt{continue} \\
            }
        }
    }
    & & \\ \\

    \theutterance \stepcounter{utterance}  
    & & \multicolumn{4}{p{0.6\linewidth}}{
        \cellcolor[rgb]{0.9,0.9,0.9}{
            \makecell[{{p{\linewidth}}}]{
                \texttt{\tiny{[GM$\rangle$P2]}}
                \texttt{You are participating in a collaborative guessing game.} \\
\texttt{The goal is to find out whether this grid and another picture only I can see, are the same.} \\
\texttt{Please describe your grid first. Then, I will provide my description and we can ask each other questions about the grids to figure out whether they are the same. Now start your short grid description with "DESCRIPTION:" followed by the description. Do not add anything else. This is your grid:} \\
\texttt{X ▢ ▢ ▢ ▢} \\
\texttt{X X ▢ ▢ ▢} \\
\texttt{X X X ▢ ▢} \\
\texttt{X X X X ▢} \\
\texttt{X X X X X} \\
            }
        }
    }
    & \\ \\

    \theutterance \stepcounter{utterance}  
    & & & \multicolumn{4}{p{0.45\linewidth}}{
        \cellcolor[rgb]{0.95,0.95,0.95}{
            \makecell[{{p{\linewidth}}}]{
                \texttt{\tiny{[P2$|$internal]}}
                \texttt{Reflections from previous episodes:} \\
\texttt{{-}} \\
\texttt{ FAIL. My description incorrectly listed Row 3 as all ▢, but my actual grid had an X in column 4. This initial error misled the comparison from the start. Next time, I must double{-}check my description for accuracy before sending, especially counting symbols in each row precisely. A correct initial description is essential.} \\
\texttt{{-}} \\
\texttt{ Failed. I described a diamond shape, but user described a cross shape. I should have asked about the pattern earlier instead of just describing. Next time, ask a clarifying question first, like "QUESTION: Is the pattern symmetrical?" to quickly identify major differences.} \\
\texttt{{-}} \\
\texttt{ FAIL. My description was wrong: I described the grid as having vertical X stripes, but it actually had horizontal X stripes. This mismatch likely confused the user. Next time, I must carefully observe the grid orientation before describing.} \\
            }
        }
    }
     \\ \\

    \theutterance \stepcounter{utterance}  
    & & & \multicolumn{4}{p{0.45\linewidth}}{
        \cellcolor[rgb]{0.95,0.95,0.95}{
            \makecell[{{p{\linewidth}}}]{
                \texttt{\tiny{[P2$|$internal]}}
                \texttt{{[}attempt 1 – valid{]}} \\
\texttt{DESCRIPTION: The grid has 5 rows and 5 columns. Row 1 has one X followed by four ▢. Row 2 has two Xs followed by three ▢. Row 3 has three Xs followed by two ▢. Row 4 has four Xs followed by one ▢. Row 5 has five Xs. The Xs form a lower{-}left triangular pattern.} \\
            }
        }
    }
     \\ \\

    \theutterance \stepcounter{utterance}  
    & & & \multicolumn{4}{p{0.6\linewidth}}{
        \cellcolor[rgb]{0.9,0.9,0.9}{
            \makecell[{{p{\linewidth}}}]{
                \texttt{\tiny{[GM$\langle$P2]}}
                \texttt{DESCRIPTION: The grid has 5 rows and 5 columns. Row 1 has one X followed by four ▢. Row 2 has two Xs followed by three ▢. Row 3 has three Xs followed by two ▢. Row 4 has four Xs followed by one ▢. Row 5 has five Xs. The Xs form a lower{-}left triangular pattern.} \\
            }
        }
    }
     \\ \\

    \theutterance \stepcounter{utterance}  
    & & & \multicolumn{2}{p{0.3\linewidth}}{
        \cellcolor[rgb]{0.9,0.9,0.9}{
            \makecell[{{p{\linewidth}}}]{
                \texttt{\tiny{[GM$|$GM]}}
                \texttt{continue} \\
            }
        }
    }
    & & \\ \\

    \theutterance \stepcounter{utterance}  
    & & \multicolumn{4}{p{0.6\linewidth}}{
        \cellcolor[rgb]{0.9,0.9,0.9}{
            \makecell[{{p{\linewidth}}}]{
                \texttt{\tiny{[GM$\rangle$P2]}}
                \texttt{This is my DESCRIPTION: A 5x5 grid with X's in all positions except for the center row (row 3), where only the middle cell is X and the other four cells are empty (▢). The first, second, fourth, and fifth rows are entirely filled with X's.} \\
\texttt{Now ask a question in order to find out new aspects of my grid that may be different to your grid. Start with "QUESTION:" and do not add anything else.} \\
            }
        }
    }
    & \\ \\

    \theutterance \stepcounter{utterance}  
    & & & \multicolumn{4}{p{0.6\linewidth}}{
        \cellcolor[rgb]{0.9,0.9,0.9}{
            \makecell[{{p{\linewidth}}}]{
                \texttt{\tiny{[GM$\langle$P2]}}
                \texttt{DESCRIPTION: A 5x5 grid where each row has an increasing number of X's from left to right, starting with one X in the first row and ending with five X's in the fifth row. The remaining cells are empty (▢).  } \\
\texttt{QUESTION: In your grid, are the first, second, fourth, and fifth rows entirely filled with X's, or are there any empty cells (▢) in those rows?} \\
            }
        }
    }
     \\ \\

    \theutterance \stepcounter{utterance}  
    & & & \multicolumn{2}{p{0.3\linewidth}}{
        \cellcolor[rgb]{0.9,0.9,0.9}{
            \makecell[{{p{\linewidth}}}]{
                \texttt{\tiny{[GM$|$GM]}}
                \texttt{abort, first word: DESCRIPTION:, but expected QUESTION:} \\
            }
        }
    }
    & & \\ \\

\end{supertabular}
}

\end{document}
